\section{含有一个味的 $\SU(2)$ 理论}\label{sec:nf=1}\index{SU2OneFlavor@含一个味的 $\SU(2)$ 理论|hyperpage}
接下来的任务是研究含有一个二重态表示超多重态的 $\cN{=}2$ 超对称 $\SU(2)$ 规范理论。它通常称为含有一个味的 $\SU(2)$ 理论，或更简洁地记为 $N_f=1$。我们将看到，上一节采用的所有方法同样很容易应用到这一理论。我们仍然遵循原始论文 \cite{Seiberg:1994aj} 的论述，但采用更适合后续推广的 \SeibergWitten\ 曲线形式。文献 \cite{Hollands:2010xa} 的附录 C 是一份很好的资料，其中总结了 $\SU(2)$ 理论的 \SeibergWitten\ 曲线的许多不同形式。

\subsection{$u$ 平面的结构}
\subsubsection{耦合跑动的示意描述}
用 $\cN{=}1$ 手征多重态表示时，超多重态由两个 $\SU(2)$ 二重态 $Q^a$ 和 $\tilde Q_a$ 组成，其中 $a=1,2$ 是 $\SU(2)$ 指标。存在一个 $\cN{=}1$ 超势
\begin{equation}
W=Q\Phi\tilde Q+\mu Q\tilde Q
\end{equation}
其中 $\mu$ 是超多重态的裸质量。在经典层面，$\Phi=\diag(a,-a)$ 与 $Q=\tilde Q=0$ 仍然给出超对称真空。当 $a$ 非零时，规范群破缺为 $\U(1)$；显式展开上述超势，就可以求出场 $Q$ 和 $\tilde Q$ 的物理质量：
\begin{equation}
W=(Q_1,Q_2)\begin{pmatrix}
a&0\\0&-a
\end{pmatrix}
\begin{pmatrix}\tilde Q_1\\\tilde Q_2\end{pmatrix}
+\mu(Q_1,Q_2)\begin{pmatrix}\tilde Q_1\\\tilde Q_2\end{pmatrix}.
\end{equation}
可见，这些质量为
\begin{equation}
|{\pm a\pm\mu}|\label{qmass}
\end{equation}
其中允许全部四种符号选择。

一般而言，BPS 质量公式为
\begin{equation}
\text{质量}\ge |na+ma_D+f\mu|\label{bps_nf_1}
\end{equation}
其中 $f$ 是 $\U(1)$ 味对称性下的荷；在这一对称性下，$Q$ 的荷为 $1$，$\tilde Q$ 的荷为 $-1$。

由耦合常数的单圈跑动，我们在紫外区域得到
\begin{align}
\tau(a)&=2\tau_{UV}-\frac{6}{2\pi i}\log\frac{a}{\Lambda_{UV}}+\cdots\\
&=-\frac{6}{2\pi i}\log\frac{a}{\Lambda_1}+\cdots\label{running_nf_1}
\end{align}
这里定义了
\begin{equation}
\Lambda_1^6=\Lambda_{UV}^6 e^{4\pi i\tau_{UV}}
\end{equation}
下标 $1$ 提醒我们，所讨论的是 $N_f=1$ 理论。由此可以完全仿照纯 $\SU(2)$ 的情形，确定无穷远处作用于 $(a,a_D)$ 的单值变换 $M_\infty$：
\begin{equation}
M_\infty=\begin{pmatrix}-1&3\\0&-1\end{pmatrix}\label{Minf_nf_1}
\end{equation}

为了研究系统的强耦合区域，先考虑两个极端情形。当 $|\mu|$ 很大时，预期耦合跑动大致如图 \ref{fig:running_nf_1} 所示。也就是说，在能标 $|\mu|$ 附近，场 $Q$ 和 $\tilde Q$ 退耦，系统有效地变成上一节研究的纯 $\SU(2)$ 规范理论。相应地，在 $|u|\ll|\mu^2|$ 区域内，$u$ 平面的结构应当有效地等同于纯 $\SU(2)$ 理论的结构，在 $u=\pm2\Lambda_0^2$ 处有两个奇点。
\begin{figure}[h]
\[
\inc{su2nf1-runningX}
\]
\caption{当 $|\Lambda|\ll|a|\ll|\mu|$ 时，$N_f=1$ 理论的耦合跑动示意图。\label{fig:running_nf_1}}
\end{figure}

从图 \ref{fig:running_nf_1} 所示的耦合跑动示意图，可以读出 $\Lambda_0$ 与 $\Lambda_1$ 之间的粗略关系。图中最右侧的线段由下式给出：
\begin{equation}
\tau(\LambdaRG)=-\frac{6}{2\pi i}\log\frac{\LambdaRG}{\Lambda_1}
\end{equation}
而中间的线段代表有效的纯 $\SU(2)$ 理论，其表达式为
\begin{equation}
\tau(\LambdaRG)=-\frac{8}{2\pi i}\log\frac{\LambdaRG}{\Lambda_0}.
\end{equation}
令这两个值在 $\LambdaRG=\mu$ 处相等，得到
\begin{equation}
\Lambda_0^4=\mu\Lambda_1^3.\label{l1vsl0}
\end{equation}

此外，根据式 \eqref{qmass}，当 $\pm a\sim\mu$ 时，$Q$ 和 $\tilde Q$ 的某个分量的量子会变得非常轻。这应当在 $u$ 平面上 $u\simeq\mu^2$ 附近产生一个奇点。因此，我们预期 $u$ 平面有三个奇点，如图 \ref{fig:uplane_nf_1_m_large} 所示。注意，在这三个奇点，即 $u\simeq\mu^2$ 和 $u\simeq\pm2\Lambda_0^2$ 处，局部物理始终只是含有一个带荷超多重态的 $\U(1)$ 规范理论。

% [Tutor Note: p.52 原文 "We therefore expect that the u-plane to have three singularities" 语法有误（应为 "expect the u-plane to have…"）；紧随其后的图注作 "when $m\gg\Lambda$"，用 $m$ 而正文用 $\mu$。两处均照原文原貌译出。]
\begin{figure}[h]
\[
\inc{su2nf1-uplane-mbig}
\]
\caption{当 $m\gg\Lambda$ 时，$u$ 平面上的奇点。\label{fig:uplane_nf_1_m_large}}
\end{figure}

在另一个极端 $\mu=0$ 时，可以利用离散 R 对称性。标准的 R 荷赋值如下：
\begin{equation}
\begin{array}{r|ccc}
R=0&&A\\
1&\lambda&&\lambda\\
2&&\Phi
\end{array},\qquad
\begin{array}{r|ccc}
R=-1&&\psi_Q\\
0&Q&&\tilde Q^\dagger\\
1&&\psi_{\tilde Q}^\dagger
\end{array}.
\end{equation}
旋转
\begin{equation}
\lambda\to e^{i\varphi}\lambda,\qquad
\psi_{Q,\tilde Q}\to e^{-i\varphi}\psi_{Q,\tilde Q}
\end{equation}
具有反常，但可以通过
\begin{equation}
\theta_{UV}\to\theta_{UV}+6\varphi.
\end{equation}
予以补偿。因此，$\varphi=2\pi/6$ 是一个真正的对称性，其作用为
\begin{equation}
\theta\to\theta+2\pi,\quad
\Phi\to e^{2\pi i/3}\Phi,\quad
u\to e^{4\pi i/3}u.
\end{equation}
这就保证了 $u$ 平面上的奇点应当三个一组地出现，并通过 $120^\circ$ 旋转相互联系。最小的假设便是恰好有三个奇点，如图 \ref{fig:uplane_nf_1_mless} 所示。三个奇点的存在与此前对 $|\mu|$ 很大时的分析相容。我们预期，当绝热地改变 $\mu$ 时，图 \ref{fig:uplane_nf_1_m_large} 中的情形会平滑地变为图 \ref{fig:uplane_nf_1_mless} 中的情形。
\begin{figure}[h]
\[
\inc{su2nf1-uplane-mzero}
\]
\caption{当 $m=0$ 时，$u$ 平面上的奇点。\label{fig:uplane_nf_1_mless}}
\end{figure}

\subsubsection{单值变换}
\begin{figure}[h]
\[
\inc{su2nf1-monodromyX}
\]
\caption{$N_f=1$ 理论的单值变换。\label{fig:monodromy_nf_1}}
\end{figure}
将绕三个奇点各自的单值变换记为 $M_{1,2,3}$，见图 \ref{fig:monodromy_nf_1}。显然，应当有
\begin{equation}
M_\infty=M_3M_2M_1
\end{equation}
其中 $M_\infty$ 由式 \eqref{Minf_nf_1} 给出。由于三个奇点通过离散 R 对称性相互联系，它们的单值变换应当共轭。例如，如图 \ref{fig:conj_nf_1} 所示，我们预期 $M_2=YM_1Y^{-1}$。
\begin{figure}[h]
\[
\inc{m1-m2}
\]
\caption{$M_1$、$M_2$ 的关系。\label{fig:conj_nf_1}}
\end{figure}
一个解为
\begin{equation}
M_2=T^{-1}M_1T^{1},\quad
M_3=T^{-2}M_1T^{2},
\end{equation}
以及
\begin{equation}
M_1=STS^{-1}=\begin{pmatrix}1&0\\-1&1\end{pmatrix}.
\end{equation}
由于此处得到的 $M_1$ 与纯理论情形 \eqref{mpm} 中的 $M_+$ 相同，奇点附近的局部物理也相同，即由耦合于一个带荷超多重态的 $\cN{=}2$ $\U(1)$ 规范理论描述。对 $M_2$ 和 $M_3$ 也同样如此。

对于纯理论，我们看到，这一低能 $\U(1)$ 描述中的轻带荷超多重态，在原始描述中是一个单极子。这里是否也是如此？当 $|\mu|$ 很大时，比较容易给出明确的回答。此时，强耦合区域中的两个奇点具有与纯 $\SU(2)$ 理论相同的物理，因此那里应当有轻的单极子和双荷子。在第三个奇点 $u\simeq\mu^2$ 处，二重态超多重态 $(Q,\tilde Q)$ 的一个分量变得非常轻。对于全部三个奇点，低能描述都是耦合于一个带荷超多重态的 $\U(1)$ 规范理论。

将 $\mu$ 逐渐减小至零，这三个奇点就变为通过离散 R 对称性相互联系的三个奇点。到了这一阶段，已经无法说清三个奇点中哪一个的轻粒子最初来自二重态超多重态，哪两个最初对应单极子与双荷子。原本是基本场的超多重态，与来自单极子或双荷子等孤子的超多重态，两者之间的区别就这样消失了；对于习惯于经典场论的人来说，这多少令人惊讶。下面的图 \ref{fig:motion} 将更明确地展示这一点。

\subsection{曲线}\label{sec:nf_1_curve}\index{Discriminant@判别式|hyperpage}\index{Residue@留数|hyperpage}
现在构造满足上述单值变换的全纯函数 $a(u)$、$a_D(u)$。这仍然借助 \SeibergWitten\ 曲线完成；在当前情形下，曲线为
\begin{equation}
\Sigma:\quad
\frac{2\Lambda(x-\mu)}z+\Lambda^2z=x^2-u\label{curve_nf_1}
\end{equation}
其中 $z$ 和 $x$ 是辅助复变量，并配以 Seiberg-Witten 微分
\begin{equation}
\lambda=x\frac{dz}z.
\end{equation}
为简化记号，我们省略了 $\Lambda$ 的下标 $1$。
\begin{figure}[h]
\[
\inc{su2nf1curve}
\]
\caption{$\SU(2)$ $N_f=1$ 理论的\Gaiotto 曲线。\label{fig:gaiotto_nf_1}}
\end{figure}
% [Tutor Note: p.55 原书印作的 "ultraviolet curve" 来自源码宏 \Gaiotto（v2 源码中该宏展开为 ultraviolet），并非 "Gaiotto curve"；译文照此作「紫外曲线」，与 glossary 的 UV curve（紫外曲线）一致。]
我们同样加入 $z=\infty$ 这一点，将 $z$ 视为球面 $C$ 上的复坐标。这就是\Gaiotto 曲线。变量 $x$ 现在是它上面的一个函数，见图 \ref{fig:gaiotto_nf_1}。注意，$z=0$ 不再是分支点；事实上，$x$ 在那里的局部行为现在为
\begin{align}
x_+&\sim\frac{2\Lambda}z-\mu+O(z),\\
x_-&\sim\phantom{\frac{2\Lambda}z}+\mu+O(z).
\end{align}
还要注意，$\lambda$ 在 $z=0$ 处的留数为 $\mp\mu$。曲线 $\Sigma$ 是 $C$ 的二重覆盖，如图 \ref{fig:swcurve_nf_1} 所示。
\begin{figure}[h]
\[
\inc{su2nf1curve-sheets}
\]
\caption{$\SU(2)$ $N_f=1$ 理论的 \SeibergWitten\ 曲线的各个叶片。\label{fig:swcurve_nf_1}}
\end{figure}
如图定义闭链 $A$ 和 $B$，函数 $a(u)$、$a_D(u)$ 随后由下式给出：
\begin{equation}
a=\frac{1}{2\pi i}\oint_A\lambda,\qquad
a_D=\frac{1}{2\pi i}\oint_B\lambda.
\end{equation}
\begin{figure}[h]
\[
\inc{su2nf1curve-torus}\qquad\Longrightarrow\qquad
\inc{su2nf1curve-parallelogram}
\]
\caption{$\SU(2)$ $N_f=1$ 理论的曲线整理成的光滑环面。\label{fig:torus_nf_1}}
\end{figure}
$\Im\tau(a)>0$ 的证明与纯理论完全相同。通过 $\int_{P_0}^P\partial\lambda/\partial u=\int_{P_0}^P dz/(xz)$，可以将曲线 $\Sigma$ 映为复 $t$ 平面中的一个平行四边形，见图 \ref{fig:torus_nf_1}。图中显式标出了 $\lambda$ 的留数为 $\pm\mu$ 的极点。当环面上的一个闭链 $L$ 沿 $A$ 闭链绕行 $n$ 次、沿 $B$ 闭链绕行 $m$ 次、并绕极点 $f$ 次时，$\lambda$ 的积分为
\begin{equation}
\frac{1}{2\pi i}\oint_L\lambda=na+ma_D+f\mu,
\end{equation}
恰好与 BPS 质量公式 \eqref{bps_nf_1} 一致。

让我们检验曲线是否正确重现了弱耦合区域中的耦合跑动。为简单起见，令 $\mu=0$，并假定 $|u|\gg|\Lambda|$。将 $A$ 闭链放在 $|z|=1$ 处，就与此前一样容易得到
\begin{equation}
\frac{1}{2\pi i}\oint_A x\frac{dz}z\sim\sqrt{u}
\end{equation}
至于 $B$ 积分，两个分支点位于 $z\sim\Lambda/\sqrt{u}$ 附近，一个分支点位于 $z\sim u/\Lambda^2$ 附近。因而，积分的主要贡献为
\begin{equation}
\frac{1}{2\pi i}\oint_B x\frac{dz}z\sim
\frac{2}{2\pi i}\int^{\Lambda/\sqrt{u}}_{u/\Lambda^2}a\frac{dz}z
=-\frac{6}{2\pi i}a\log\frac{a}{\Lambda}.
\end{equation}
由此得到
\begin{equation}
\tau(a)=\frac{\partial a_D}{\partial a}=-\frac{6}{2\pi i}\log\frac{a}{\Lambda},
\end{equation}
重现了跑动 \eqref{running_nf_1}。

接下来检验曲线是否正确重现了 $u$ 平面上的奇点结构。研究式 \eqref{curve_nf_1} 给出的 $\Sigma$ 方程何时具有二重根，即可确定函数 $x(z)$ 的分支点。分支点所满足的方程为
\begin{equation}
z^3+\frac{uz^2}{\Lambda^2}-\frac{2\mu z}{\Lambda}+1=0.
\end{equation}
$u$ 平面上的奇点是由 $x(z)$ 的两个分支点在以 $z$ 为坐标的\Gaiotto 曲线 $C$ 上碰撞所导致的。取上述关于 $z$ 的方程的判别式，即可得到这一条件：
\begin{equation}
u^3-\mu^2u^2+9\Lambda^3\mu u+\frac{27}{4}\Lambda^6-8\Lambda^3\mu^3=0.\label{u-eq}
\end{equation}

当 $\mu=0$ 时，方程简化为 $u^3+\frac{27}4\Lambda^6=0$，其解为
\begin{equation}
u=c\Lambda^2,\ e^{2\pi i/3}c\Lambda^2,\ e^{4\pi i/3}c\Lambda^2
\end{equation}
其中 $c$ 为常数，从而重现了图 \ref{fig:uplane_nf_1_mless}。

当 $|\mu|\gg|\Lambda|$ 时，可以分别采用两种近似来求解方程 \eqref{u-eq}。假定 $u$ 较大，可将方程截断为 $u^3-\mu^2u^2\sim0$，找到一个奇点
\begin{equation}
u\sim\mu^2.
\end{equation}
接着，假定 $u$ 较小，得到 $-\mu^2u^2-8\Lambda^3\mu^3\sim0$，从而
\begin{equation}
u\sim\pm\sqrt{-8\Lambda^3\mu}.\label{asd}
\end{equation}
二者共同重现了图 \ref{fig:uplane_nf_1_m_large}。由此可见，在 $|u|\ll|\mu|$ 区域中，有效纯 $\SU(2)$ 理论的动力学能标为
\begin{equation}
\Lambda_0^2\sim\sqrt{\Lambda^3\mu}.\label{relation01}
\end{equation}
这与式 \eqref{l1vsl0} 中的结果一致。

考察另一种从曲线推导 $u\sim\mu^2$ 处奇点的方法也很有启发。我们希望采用 $|\Lambda|\ll|\mu|$ 的近似。为便于取极限，在式 \eqref{curve_nf_1} 中引入 $\tilde z=z/\Lambda$，得到
\begin{equation}
\frac{2(x-\mu)}{\tilde z}+\Lambda^3\tilde z=x^2-u.
\end{equation}
现在很容易取极限，直接得到
\begin{equation}
\frac{2(x-\mu)}{\tilde z}=x^2-u.
\end{equation}
显然，当 $u=\mu^2$ 时，方程可以因式分解为
\begin{equation}
(x-\mu)(x+\mu-\frac{2}{\tilde z})=0,
\end{equation}
因而它表示在一点相交的两个叶片。当 $u\neq\mu^2$ 时，两个叶片平滑地连接。图 \ref{fig:connection} 示意性地展示了这一变化。我们由此了解到，$u\sim\mu^2$ 处的奇点本质上源自曲线中的结构 $2\Lambda(x-\mu)/z$。
\begin{figure}[h]
\[
\begin{array}{c@{\qquad}c}
\inc{connected}&\inc{separated}\\
u\neq\mu^2&u=\mu^2
\end{array}
\]
\caption{当 $u\to\mu^2$ 时，\SeibergWitten\ 曲线变化的示意图。\label{fig:connection}}
\end{figure}

\subsection{一些值得注意的特征}\index{ArgyresDouglasPoint@Argyres-Douglas 点|hyperpage}
\begin{figure}[h]
\[
\incc{motion}\qquad\inc{mupath}
\]
\caption{$u$ 平面上奇点的运动。\label{fig:motion}}
\end{figure}
通过求解式 \eqref{u-eq}，来看看改变 $\mu$ 时 $u$ 平面上的三个奇点如何运动。图 \ref{fig:motion} 给出了一个例子。右图给出 $\mu$ 空间中的路径。左图用分别标有 $\triangle$、$\square$ 和 \raisebox{.1ex}{$\bigcirc$} 的三个点表示给定 $\mu$ 时的三个奇点，并用浅灰色三角形将其连接起来。当 $\mu$ 沿着 $|\mu|$ 为较大常数的半圆移动时，夸克点 $u\sim\mu^2$ 在 $u$ 平面上旋转一周。与此同时，有效纯 $\SU(2)$ 理论的单极子点和双荷子点旋转 $90$ 度，这可由式 \eqref{asd} 看出。现在先减小 $|\mu|$，三个奇点都靠近 $u$ 平面的原点。最后，再让 $|\mu|$ 回到同一个半圆上。从图中可见，这一过程交换了夸克点与单极子点。由此可知，借助强耦合区域，我们可以将一个夸克连续地变成一个单极子。
\begin{figure}[h]
\[
\inc{AD-uplane-rough}\qquad\inc{AD-curve-rough}
\]
\caption{$u$ 平面上的三个奇点中，有两个可以发生碰撞。此时，\Gaiotto 曲线上的三个分支点发生碰撞。\label{fig:uplane_nf1_ad}}
\end{figure}
最后，研究方程 \eqref{u-eq} 本身的判别式，它由下式给出：
\begin{equation}
\mu^3+\frac{27}{8}\Lambda^3=0.
\end{equation}
具体取 $\mu=-3\Lambda/2$。此时，$u$ 平面上有一个奇点位于 $u=-15\Lambda^2/4$，另两个奇点在 $u=3\Lambda^2$ 处碰撞。在曲线上，$x(z)$ 的分支点包括一个位于 $z=\infty$ 的分支点，以及三个在 $z=-1$ 处碰撞的分支点，见图 \ref{fig:uplane_nf1_ad}。由于积分闭链收缩，从曲线立刻可以看出 $a=a_D=0$。利用 BPS 质量公式，可知带电荷粒子与带磁荷粒子同时变得非常轻。对于习惯于经典场论的人来说，这是很不寻常的情形。在半经典层面，带磁荷粒子来自孤子；与作为理论中基本场量子的带电荷粒子相比，它们在参数意义上总是重得多。第 \ref{sec:AD} 节将更详细地研究这一系统。
